\chapter{Metodología}
\label{cap:metodologia}

Este capítulo presenta la metodología propia derivada de la OWASP WSTG,
descrita en el capítulo~\ref{cap:marco-teorico}, así como el proceso
concreto seguido para aplicarla sobre \breachlab.

\section{Derivación de una metodología ágil a partir de la OWASP WSTG}
\label{sec:derivacion-metodologia}

La OWASP WSTG v4.2 organiza sus pruebas en once categorías (gestión de la
configuración, gestión de identidad, autenticación, autorización, gestión de
sesión, validación de entradas, manejo de errores, criptografía, lógica de
negocio, del lado del cliente y pruebas de API), con más de noventa pruebas
individuales en total. Recorrer la guía completa exige un equipo con
disponibilidad y conocimiento amplios, algo razonable en una auditoría
externa de alcance general, pero difícilmente asumible por un equipo de
desarrollo pequeño o una pyme que necesite validar de forma recurrente los
riesgos más críticos de su propia aplicación sin subcontratar una auditoría
completa.

A partir de esta observación se deriva la metodología propia aplicada en
este trabajo, que agiliza la WSTG mediante dos decisiones:

\begin{itemize}
\item \textbf{Acotar el alcance a un subconjunto priorizado de pruebas.} En
  lugar de recorrer las once categorías de la WSTG, se seleccionan
  únicamente las pruebas correspondientes a las categorías del OWASP Top~10
  con mayor prevalencia e impacto documentado~\citelib{OWASP21}
  (inyección, XSS, CSRF, autenticación, control de acceso y configuración de
  seguridad), descartando de forma explícita categorías de menor
  probabilidad en el contexto de una aplicación pequeña (por ejemplo,
  pruebas de criptografía avanzada o de lógica de negocio compleja), que
  quedan como trabajo futuro (apartado~\ref{sec:trabajo-futuro}).
\item \textbf{Consolidar las fases por vulnerabilidad en un único ciclo
  corto.} Frente al desglose detallado de sub-pasos de cada prueba WSTG, se
  define un ciclo único de seis pasos (apartado~\ref{sec:proceso}),
  reconocimiento, identificación, explotación, evidencia, verificación de la
  mitigación y valoración de impacto, aplicado de forma idéntica a cada
  vulnerabilidad, lo que reduce el tiempo de preparación por prueba y
  simplifica la comparación de resultados entre vulnerabilidades distintas.
\end{itemize}

\needspace{15\baselineskip}
El resultado es una metodología más rápida de aplicar y más accesible para
un equipo con recursos limitados, a costa de una cobertura menor que la de
la WSTG completa. El capítulo~\ref{cap:comparativa} cuantifica esta
diferencia de cobertura y discute en qué contextos resulta preferible cada
una de las dos aproximaciones.

\section{Proceso aplicado}
\label{sec:proceso}

Para cada una de las seis vulnerabilidades introducidas deliberadamente en \breachlab se ha
seguido el mismo proceso de trabajo, descrito a continuación, con el fin de
que los resultados del capítulo~\ref{cap:explotacion} sean comparables entre
sí y reproducibles.

\begin{enumerate}
\item \textbf{Reconocimiento:} identificación del endpoint o vista implicada
  y de la entrada de usuario que llega hasta él, revisando el código fuente
  de \texttt{backend/api.py} y de las páginas Astro correspondientes.
\item \textbf{Identificación:} localización exacta de la línea de código
  donde el dato no confiable se usa de forma insegura (concatenación en una
  consulta SQL, asignación a \verb|innerHTML|, ausencia de comprobación de
  propietario, etc.).
\item \textbf{Explotación:} construcción de un payload o petición que
  demuestre el fallo en la rama vulnerable (por defecto, sin
  \texttt{?safe=1}), utilizando el navegador, \texttt{curl} o herramientas
  especializadas según el caso.
\item \textbf{Evidencia:} captura del resultado de la explotación junto con
  la línea correspondiente registrada en \texttt{backend/logs/attacks.log}.
\item \textbf{Verificación de la mitigación:} repetición de la misma
  petición con \texttt{?safe=1} para comprobar que la contramedida
  implementada bloquea el ataque, y comparación del fragmento de código
  vulnerable frente al seguro.
\item \textbf{Valoración de impacto:} cálculo del vector y puntuación
  CVSS v3.1 correspondientes (véase el apartado~\ref{sec:cvss}).
\end{enumerate}

\needspace{15\baselineskip}
\section{Correspondencia con la OWASP WSTG}

Aunque el proceso de seis pasos anterior es común a las seis
vulnerabilidades, cada una de ellas se corresponde además con una prueba
concreta del catálogo de la OWASP Web Security Testing Guide
(WSTG)~\citelink{OWASPWSTG} mencionada en el
capítulo~\ref{cap:marco-teorico}, lo que permite situar el trabajo dentro de
un marco de referencia reconocido en la industria en lugar de depender únicamente de una
metodología ad-hoc. La tabla~\ref{tab:wstg-mapeo} recoge todas las
correspondencias entre las vulnerabilidades y sus identificadores dentro de WSTG.

\begin{table}[htbp]
  \centering
  \small
  \begin{tabular}{lll}
    \toprule
    \tabhead{Vulnerabilidad en \breachlab} & \tabhead{Identificador WSTG} & \tabhead{Prueba} \\
    \midrule
    Inyección SQL                    & WSTG-INPV-05 & Testing for SQL Injection \\
    XSS basado en DOM                & WSTG-CLNT-01 & Testing for DOM-based XSS \\
    CSRF                             & WSTG-SESS-05 & Testing for CSRF \\
    Enumeración de usuarios          & WSTG-IDNT-04 & Testing for Account Enumeration \\
    Fuerza bruta / bloqueo de cuenta & WSTG-ATHN-03 & Testing for Weak Lock Out Mechanism \\
    IDOR                             & WSTG-ATHZ-04 & Testing for IDOR \\
    Config.\ insegura (clickjacking) & WSTG-CLNT-09 & Testing for Clickjacking \\
    \bottomrule
  \end{tabular}
  \caption{Correspondencia entre las vulnerabilidades de \breachlab y las
    pruebas de la OWASP WSTG v4.2.}
  \label{tab:wstg-mapeo}
\end{table}

\section{Herramientas empleadas}

Además del navegador web y de las herramientas de desarrollador integradas
en él, se han utilizado: \texttt{curl}, para reproducir peticiones HTTP de
forma exacta y guardarlas como evidencia; \texttt{sqlmap}~\citelink{SQLMAP},
para automatizar
la detección y explotación de la inyección SQL; y el propio registro
\texttt{backend/logs/attacks.log}, seguido en tiempo real con
\verb|tail -f|, como fuente de evidencia objetiva independiente de las
capturas de pantalla.

\section{Criterios de severidad}

La severidad de cada vulnerabilidad se expresa mediante su vector y
puntuación CVSS v3.1, calculados con la calculadora oficial de
FIRST.org~\citelink{CVSSCALC} a
partir del vector razonado en el apartado correspondiente del
capítulo~\ref{cap:explotacion}.
Cuando una vulnerabilidad no es explotable de forma aislada, como sería el caso de la
configuración de seguridad insegura del apartado~\ref{sec:secmisconfig}, se documenta como un
control compensatorio ausente que amplifica el impacto de otra
vulnerabilidad, en lugar de forzar una puntuación CVSS independiente que no
reflejaría correctamente la naturaleza de dicha vulnerabilidad.
